\documentclass[10pt,a4paper]{amsart}
\usepackage[margin=19mm]{geometry}
\usepackage{amsmath,amssymb,mathtools}
\usepackage[scr=rsfs,cal=euler]{mathalfa}
\usepackage{xfrac}
\usepackage[unicode,hidelinks,bookmarks=false]{hyperref}
\newcommand{\ie}{i.e.\ }
\newcommand{\cf}{cf.\ }
\newcommand{\resp}{respectively\ }
\newcommand{\Subsection}{\subsection}
\providecommand{\nobreakdash}{\nobreak\hbox{-}}
\setlength{\emergencystretch}{2em}
\multlinegap0pt
\usepackage{bm}
\usepackage{bbm}
\usepackage{dsfont}
\usepackage{xspace}
\usepackage{cancel}
\usepackage{mathtools}
\usepackage{amsthm}
\makeatletter
\@ifundefined{theorem}{\newtheorem{theorem}{Theorem}}{}
\makeatother
\title[The chromatic number of Euclidean space with dense color cla]{The chromatic number of Euclidean space with dense color classes}
\author{Maxim Didin, Vsevolod Voronov}
\date{Source: arXiv:2607.19946v1 (CC BY 4.0)}
\begin{document}
\maketitle

\noindent\emph{Source and licence.} The chromatic number of Euclidean space with dense color classes. Version arXiv:2607.19946v1 (CC BY 4.0), distributed under the Creative Commons Attribution 4.0 International licence, as recorded in the arXiv licence metadata for this exact version and shown on the abstract page https://arxiv.org/abs/2607.19946v1. Full rights notice: licenses/arxiv-2607.19946-CC-BY-4.0.txt.\par\smallskip

\noindent\textit{Editorial scope.} Counted unit: the Theorem whose source label is \texttt{thm:main} in the article (source0.tex lines 148--155, source label \texttt{thm:main}), delivered together with the article's own complete original proof of it (source0.tex lines 157--246, one \texttt{proof} environment of 90 lines). No mathematics is written, repaired or re-derived: statement and proof are the article's own text line for line. Required context retained uncounted: none. Every definition, display and label of those lines is retained unchanged. Omitted from this package and not redistributed: 1--147, 156--156, 247--579. Nothing inside a retained excerpt is omitted or altered: every retained body is delivered exactly as the article's lines are written, so no omission receipt of any kind accompanies this package. Citations: the retained excerpts carry no citation command at all, so no bibliography entry of the article is transcribed and no reference is redistributed. Reference adaptation: none was needed and none was made; every reference inside the delivered unit resolves inside it, so no unresolved reference or citation is produced. Printed slips kept literal: no printed word, symbol, hypothesis or number of the counted statement or of its proof was altered. Layout and class: the article's own class is \texttt{article} with packages including inputenc, fontenc, babel, amscd, amssymb, amsthm, amsfonts, amsmath, hyperref, graphicx; the wrapper is the lane's pinned \texttt{amsart} editorial wrapper at 10pt on A4, which restates the theorem-like environments the retained text needs and adds the article's own macro definitions that the retained text needs (none), with extra wrapper packages (bm, bbm, dsfont, xspace, cancel, mathtools). The delivered numbering is the unit's own. Assets: no retained span contains a figure environment, a graphics-inclusion command, a PDF-page inclusion, a table or a TikZ picture; no image file is copied into this package, the package declares no asset dependency and no image file is redistributed with it. Licence: version arXiv:2607.19946v1 of this article is distributed under the Creative Commons Attribution 4.0 International licence, as recorded in the arXiv licence metadata for this exact version and shown on the abstract page https://arxiv.org/abs/2607.19946v1. A full rights notice is delivered at licenses/arxiv-2607.19946-CC-BY-4.0.txt. Characters: every retained excerpt is pure ASCII, so the delivered engine, XeLaTeX with its default Latin Modern text font, reports no missing character.
\begin{theorem}
Suppose that $\mathbb{R}^n$ admits a proper coloring with $k$ colors.
Then
\[
    \chi_{ed}(\mathbb{R}^n)\leq nk+1.
\]
\label{thm:main}
\end{theorem}

\begin{proof}
We shall say that the coordinates of the points of a set
$X\subset\mathbb{R}^n$ are algebraically independent if the union of the
sets of first coordinates, second coordinates, and so on up to the
$n$-th coordinates, in other words, the union of the projections of $X$ onto
the coordinate axes is algebraically independent over $\mathbb{Q}$.

Let $X$ be an everywhere dense subset of $\mathbb{R}^n$ whose coordinates
are algebraically independent over $\mathbb{Q}$. Observe that such a set exists. Indeed, if $X$ is finite
or countable, then the set of such points is contained in a countable union
of proper algebraic subsets of \(\mathbb R^n\). Hence it has empty
interior. In particular, every nonempty open ball contains a point that
can be added.  In particular, for every $\varepsilon>0$,
the construction can be carried out so that $X=X(\varepsilon)$ is an
$\varepsilon$-net in $\mathbb{R}^n$.

Let $A_1,\dots,A_k$ be the color classes in the original proper coloring
of $\mathbb{R}^n$. We shall construct a proper coloring of $\mathbb{R}^n$
with $nk+1$ dense color classes.

Partition $X$ into $n$ subsets
$X_1,\dots,X_n$, each dense in $\mathbb{R}^n$. Define
\[
    A'_{i,j}
    =
    X_j\cup
    \left(
        A_i\setminus
        \bigcup_{x\in X_j}S^{n-1}_1(x)
    \right),
    \qquad
    1\leq i\leq k,\quad 1\leq j\leq n,
\]
where $S^{n-1}_1(x)$ denotes the unit sphere centered at $x$.

Let
\[
    A^*
    =
    \bigcup_{x_j\in X_j,\;1\leq j\leq n}
    \bigcap_{1\leq j\leq n} S^{n-1}_1(x_j).
\]
Thus, $A^*$ is the set of all intersections of $n$ unit spheres whose
centers are chosen from $X_1,\dots,X_n$, one center from each set.

The set $A^*$ is everywhere dense. Indeed, given any point
$y\in\mathbb{R}^n$, choose points
\[
z_1,\dots,z_n\in S^{n-1}_1(y)
\]
such that the vectors $z_1-y,\dots,z_n-y$ are linearly independent.
If the points $x_j\in X_j$ are chosen sufficiently close to $z_j$, then
the corresponding $n$ spheres have an intersection point arbitrarily
close to $y$. This follows, for example, from the implicit function
theorem. Hence every neighborhood of $y$ intersects~$A^*$.

The sets
\[
    B_{i+n(j-1)}=A'_{i,j},
    \qquad 1\leq i\leq k,\quad 1\leq j\leq n,
\]
together with
\[
    B_{nk+1}=A^*
\]
form a cover of $\mathbb{R}^n$. Indeed, any point not covered by the sets
$A'_{i,j}$ belongs, by construction, to the intersection of $n$ unit
spheres, and hence belongs to~$A^*$.

Each of the sets $A'_{i,j}$ is unit distance free: the set $A_i$ is
unit distance free by assumption, points of $X_j$ cannot be at unit
distance from one another because of the algebraic independence of the
coordinates of $X$, and the spheres removed from $A_i$ eliminate pairs
consisting of a point of $X_j$ and a point of $A_i$. Moreover, the set
$A^*$ is unit distance free, again by the algebraic independence of the
coordinates of the centers.

It remains to pass from the cover by the everywhere dense
unit distance free sets
$B_1,\dots,B_{nk+1}$ to a partition into everywhere dense sets
$C_1,\dots,C_{nk+1}$. For example, choose countable, pairwise disjoint,
everywhere dense subsets
\[
    B'_i\subset B_i,\qquad 1\leq i\leq nk+1,
\]
and put $B'_i$ into $C_i$. Assign each remaining point to one of the
indices $i$ for which it belongs to $B_i$. Thus $C_i\subseteq B_i$ for
every $i$, so the resulting partition is proper, and every $C_i$ is
everywhere dense.
\end{proof}

\end{document}
